% ৩.৭ কোড (Code) - ধারণা
\section{কোড: ধারণা ও প্রয়োজন}

\begin{itemize}
\item Code কী, data ও character representation
\item Numeric ও alphanumeric code
\item কেন computer code ব্যবহার করে
\end{itemize}

\begin{learningbox}
এই টপিক শেষে শিক্ষার্থী পারবে:
\begin{itemize}
\item computer code-এর প্রয়োজনীয়তা ব্যাখ্যা করতে।
\item numeric code ও alphanumeric code আলাদা করতে।
\item character representation-এর ধারণা লিখতে।
\end{itemize}
\end{learningbox}

\begin{definitionbox}{Computer code}
কম্পিউটার কোনো character, digit বা symbol-কে নির্দিষ্ট binary pattern দিয়ে প্রকাশ করে। character বা data প্রকাশের এই নির্দিষ্ট binary নিয়মকে computer code বলা হয়।
\end{definitionbox}

\begin{keypointbox}{কেন code দরকার}
\begin{itemize}
\item computer সরাসরি বাংলা/ইংরেজি অক্ষর বোঝে না; binary pattern বোঝে।
\item text, number, symbol ও control instruction সংরক্ষণে code দরকার।
\item একই character সব system-এ একইভাবে বোঝাতে standard code দরকার।
\end{itemize}
\end{keypointbox}

\begin{center}
\fbox{\parbox{0.88\textwidth}{
\centering
\textbf{Diagram placeholder}\\[4pt]
Character ``A'' $\rightarrow$ code table lookup $\rightarrow$ binary code $\rightarrow$ memory/storage
}}
\end{center}

\begin{practicebox}
\begin{enumerate}
\item Computer code কী?
\item Numeric code ও alphanumeric code-এর পার্থক্য লেখ।
\item Unicode কেন গুরুত্বপূর্ণ?
\end{enumerate}
\end{practicebox}
